% =================
% A NOTE TO THE READER — formal content notice (non-diegetic)
% Official page before the fictional immersion begins
% =================
\cleardoublepage
\thispagestyle{empty}
\begingroup
\setlength{\fboxsep}{8mm}
\begin{center}
{\color{sangue}\scshape\Large A NOTE TO THE READER}\\[1.2em]
{\color{sangue}\rule{0.4\paperwidth}{0.6pt}}\\[2.2em]
\end{center}

\noindent This book contains, in its fictional part, direct representations of:

\begin{itemize}
\itemsep0.4em
\item institutional psychiatric violence --- including electroshock, lobotomy, and imprisonment in hospital-colonies (Barbacena Hospital-Colony, 1903--1994);
\item famine, death, and mass burial of children and adults (the Ceará droughts, 1915--1932);
\item detailed medical and necroscopic descriptions, with explicit clinical language;
\item death, grief, abandonment, and long-term isolation;
\item second-person narrative that directly summons the reader to occupy the place of \textit{Patient 1,263}, with simulated reading instructions.
\end{itemize}

\vspace{1em}

\noindent \textbf{About the simulated instructions:} no instruction found in this book produces any real physical, chemical, or biological effect. The ``contamination'' is a narrative device. The pages were printed with ordinary paper and ink, with no substance applied. The reader may stop reading at any time, at any fragment, with no consequence other than those the story itself suggests.

\vspace{1em}

\noindent \textbf{About dreams, sleep and symptoms:} there is no evidence that reading this book induces specific dreams, alters the sleep--wake cycle or produces any physical symptom. The descriptions of bodily reaction are a formal device, not an empirical claim.

\vspace{1em}

\noindent \textbf{About the clinical codes and scales:} the diagnostic protocols, ICD codes and assessment scales appearing in the fictional part are a literary construction. The disorders named are real clinical conditions, but their application to a narrative agent has no medical basis whatsoever. \textbf{This book does not diagnose the reader}: it simulates a diagnosis, as a device of fiction.

\vspace{1em}

\noindent \textbf{About the fiction/reality distinction:} the characters Joaquim Antônio Correia, Dr. Heitor Oliveira, Dr. Lúcia Mendes, and the creature are fictional. The government pens (1915--1932) and the Barbacena Hospital-Colony (1903--1994) are documented historical facts. The full distinction is set out in the \textit{Historical Note}, at the end of the volume (p. \pageref{nota-historica}).

\vspace{1em}

\noindent \textbf{If you are sensitive to any of these themes:} this work may not be suitable for you right now. There is no shame in postponing the reading or setting it aside. Books should not be endurance tests.

\vspace{1em}

\noindent \textbf{If you need support:} Centro de Valorização da Vida --- CVV, 188 (toll-free, 24 hours, Brazil) or \texttt{cvv.org.br}; 988 Suicide \& Crisis Lifeline (United States) or \texttt{988lifeline.org}; your local mental health services; or your trusted network.

\vspace{2.2em}
\begin{center}
{\color{sangue}\rule{0.4\paperwidth}{0.6pt}}\\[1em]
{\small\itshape The Archivist}
\end{center}
\endgroup
\null\vfill
\clearpage
